% TOP_PROBLEM: {"id":30007590,"problem_number":"LOCAL-30007590","title":"Scaling behavioral interventions","category_id":21,"category":{"id":21,"name":"economics","display_name":"Economics","description":"Open economic theory statements, published research agendas, and proposed scoped specifications; question provenance and historical priority are qualified per record.","slug":"economics","order_index":21,"created_at":"2026-10-08T00:00:00Z"},"status":"research_question","external_url":null,"legacy_ids":[],"published":false,"statement":"Estimate how a behavioral intervention’s effect changes when scaled across implementers, populations, and treatment saturation while accounting for selection and spillovers.","economics":{"original_id":79,"field":"Behavioral and experimental economics","kind":"empirical","formalization_basis":"adapted_research_question","source_status":"research_agenda","priority_status":"earliest_found","priority_note":"The March 2020 primary manuscript verifies a scaling question and a remaining prediction gap. Earlier scaling literature exists; this is not a verified original problem statement.","earliest_verified_reference":{"authors":["Stefano DellaVigna","Elizabeth Linos"],"title":"RCTs to Scale: Comprehensive Evidence from Two Nudge Units","year":2020,"url":"https://eml.berkeley.edu/~sdellavi/wp/NudgeToScale2020-03-20.pdf","doi":null,"locator":"Author manuscript dated March 2020, §1, printed p. 1; end of §4 and start of §5, printed p. 21 (PDF pp. 2 and 22).","evidence_summary":"The introduction frames scaling as a research question. At the end of the forecast analysis the authors say a larger sample of trials is needed to resolve whether forecasters identify more effective interventions. Average-effect comparisons alone do not identify a same-intervention causal scale effect."},"current_reference":{"authors":["Stefano DellaVigna","Elizabeth Linos"],"title":"RCTs to Scale: Comprehensive Evidence from Two Nudge Units","year":2020,"url":"https://eml.berkeley.edu/~sdellavi/wp/NudgeToScale2020-03-20.pdf","doi":null,"locator":"Author manuscript dated March 2020, §1, printed p. 1; end of §4 and start of §5, printed p. 21 (PDF pp. 2 and 22).","evidence_summary":"The introduction frames scaling as a research question. At the end of the forecast analysis the authors say a larger sample of trials is needed to resolve whether forecasters identify more effective interventions. Average-effect comparisons alone do not identify a same-intervention causal scale effect."},"formalization_note":"The source supports a scaling and forecasting agenda. This formulation extends it to a clearly defined same-intervention scale contrast; it does not recycle the paper’s answered average comparison as an open question."},"statusQualification":"Published question or research agenda verified at the cited locator. The mathematical specification is adapted for this catalogue and includes additional assumptions; source publication does not establish first-ever priority or certify this exact model remains unsolved. The source supports a scaling and forecasting agenda. This formulation extends it to a clearly defined same-intervention scale contrast; it does not recycle the paper’s answered average comparison as an open question.","statusReviewedAt":"2026-10-08"}
% FIRST_POSTED: 2026-10-08
% ENTRY_KIND: statement_only
% !TeX program = lualatex
\documentclass[11pt]{article}
\usepackage{fontspec}
\usepackage[a4paper,margin=25mm]{geometry}
\usepackage{amsmath,amssymb,mathtools}
\usepackage{xurl}
\usepackage[hidelinks]{hyperref}
\newcommand{\sref}[2]{\hyperlink{source:#1:#2}{[S#2]}}
\setlength{\emergencystretch}{3em}
\begin{document}
% problemId:30007590
\section*{Scaling behavioral interventions}\label{30007590}
\subsection{Definitions and mathematical statement}
Index interventions by \(j\) and implementation scale by \(s\), defined as the fraction of an eligible population reached. Let \(Y_{ij}(d,s)\) be participant \(i\)'s outcome under receipt \(d\in\{0,1\}\), allowing implementation capacity and spillovers to depend on scale. Define the direct treatment effect \(\tau_j(s)=\mathbb E[Y_{ij}(1,s)-Y_{ij}(0,s)]\). Let \(X_j\) contain preregistered intervention features and \(C_j(s)\) contain measured implementation features. The prediction target for a large-scale rollout is
\[\widehat\tau_j(s_H)=g\bigl(\widehat\tau_j(s_L),X_j,C_j(s_L),C_j(s_H)\bigr),\qquad R=\mathbb E[(\widehat\tau_j(s_H)-\tau_j(s_H))^2].\]
Estimate and validate \(g\) on interventions held out from training, using complete trial registries rather than only published significant effects. Separate differences caused by publication selection, population composition, intervention versions, and administrative constraints from the causal effect of changing scale itself. Observing academic and government trials at different scales does not identify \(\tau_j(s_H)-\tau_j(s_L)\) for a common intervention. Identification of that contrast requires randomized saturation, replicated implementation, or explicit transport assumptions. Assess both rank prediction and magnitude prediction, and include uncertainty from estimated trial effects in \(R\). The remaining empirical agenda concerns which characteristics provide reliable forecasts across interventions and institutions; the average academic-versus-nudge-unit comparison already has evidence.

\par\textbf{Formulation and scope.} The source supports a scaling and forecasting agenda. This formulation extends it to a clearly defined same-intervention scale contrast; it does not recycle the paper’s answered average comparison as an open question.

\par\textbf{Open-question provenance.} An explicit question or research agenda was verified at the source locator. This does not by itself certify that every later version remains unresolved \sref{30007590}{1}.

\par\textbf{Historical priority.} The March 2020 primary manuscript verifies a scaling question and a remaining prediction gap. Earlier scaling literature exists; this is not a verified original problem statement.

\subsection{Short English statement}
Estimate how a behavioral intervention’s effect changes when scaled across implementers, populations, and treatment saturation while accounting for selection and spillovers.

\subsection{Sources}
\begin{itemize}\raggedright
\item[\textnormal{[S1]}]\hypertarget{source:30007590:1}{} Stefano DellaVigna, Elizabeth Linos. 2020. RCTs to Scale: Comprehensive Evidence from Two Nudge Units. Author manuscript dated March 2020, §1, printed p. 1; end of §4 and start of §5, printed p. 21 (PDF pp. 2 and 22)..
\url{https://eml.berkeley.edu/~sdellavi/wp/NudgeToScale2020-03-20.pdf}.
{\small\textit{Earliest verified question/agenda reference; historical priority qualified below. The introduction frames scaling as a research question. At the end of the forecast analysis the authors say a larger sample of trials is needed to resolve whether forecasters identify more effective interventions. Average-effect comparisons alone do not identify a same-intervention causal scale effect.}}
\end{itemize}
\end{document}
